% !TEX root = ../main.tex
\corrigesChapitre{11}{Interactions fondamentales et champs}
\corrigesSouples

\begin{blocLivre}{cor-11-1}
\begin{corrige}{1}{Électriser une règle}
\demarche{Une charge négative signifie un excès d'électrons : $n=\dfrac{|q|}{e}$. La charge totale se conserve lors du frottement.}
\begin{questions}
\item La règle a \res{\text{gagné}} des électrons : $n=\dfrac{\num{3,2e-9}}{\num{1,60e-19}}=\res{\num{2,0e10}}$ électrons.
\item Les électrons viennent du chiffon, initialement neutre : il porte la charge $\res{\qty{+3,2}{nC}}$ (conservation de la charge).
\item Par \res{\text{influence}}, la face des morceaux de papier tournée vers la règle se charge positivement ; plus proche que la face
négative, elle est plus attirée qu'elle n'est repoussée.
\item Le métal est \res{\text{conducteur}} : les électrons gagnés s'écoulent par la main et le corps vers la Terre.
\end{questions}
\piege{Oublier de convertir : \qty{3,2}{nC} $=\qty{3,2e-9}{C}$ (nano $=10^{-9}$). Le nombre d'électrons est toujours un entier positif, très grand.}
\end{corrige}
\end{blocLivre}
\begin{blocLivre}{cor-11-2}
\begin{corrige}{2}{Appliquer la loi de Coulomb}
\demarche{Écrire la loi littérale, convertir en unités SI (C et m), calculer, puis tracer les deux forces à partir de leur point d'application.}
\begin{questions}
\item Les charges sont de signes opposés : elles \res{\text{s'attirent}}.
\item $F=k\,\dfrac{|q_\mathrm{A}\,q_\mathrm{B}|}{d^2}=\num{9,0e9}\times\dfrac{\num{2,0e-9}\times\num{3,0e-9}}{(\num{5,0e-2})^2}=\res{\qty{2,16e-5}{N}}$.
\item Les deux forces, de même valeur, sont représentées par des flèches de \qty{2,16}{cm}, chacune dirigée vers l'autre charge.
\item $F$ est proportionnelle à $\dfrac{1}{d^2}$ : si $d$ double, $F$ est \res{\text{divisée par 4}} : $\res{\qty{5,4e-6}{N}}$.
\end{questions}
\schema[corrige]{11-ex02}
\retiens{Principe des actions réciproques : $\vec{F}_{\mathrm{A}/\mathrm{B}}=-\vec{F}_{\mathrm{B}/\mathrm{A}}$. Les deux forces ont toujours la même valeur, même si
les charges sont différentes.}
\end{corrige}
\end{blocLivre}
\begin{blocLivre}{cor-11-3}
\begin{corrige}{3}{La Terre et la Lune}
\demarche{Loi de la gravitation : $F=G\,\dfrac{m_\mathrm{A}\,m_\mathrm{B}}{d^2}$, avec les distances en mètres et les masses en kilogrammes.}
\begin{questions}
\item $F_{\mathrm{T}/\mathrm{L}}=\num{6,67e-11}\times\dfrac{\num{5,97e24}\times\num{7,35e22}}{(\num{3,84e8})^2}=\res{\qty{1,98e20}{N}}$.
\item Même valeur, $\res{\qty{1,98e20}{N}}$, même direction, sens opposé (schéma ci-dessous).
\item $F_{\mathrm{S}/\mathrm{T}}=\num{6,67e-11}\times\dfrac{\num{1,99e30}\times\num{5,97e24}}{(\num{1,50e11})^2}=\res{\qty{3,52e22}{N}}$.
\item $\dfrac{\num{3,52e22}}{\num{1,98e20}}\approx\res{178}$ : le Soleil attire la Terre environ 180 fois plus que la Lune.
\end{questions}
\schema[corrige]{11-ex03-cor}
\end{corrige}
\end{blocLivre}
\begin{blocLivre}{cor-11-4}
\begin{corrige}{4}{Le champ de pesanteur terrestre}
\demarche{$g=G\,\dfrac{M_\mathrm{T}}{d^2}$, où $d$ est la distance au \emph{centre} de la Terre : $d=R_\mathrm{T}$ au sol, $d=R_\mathrm{T}+h$ en altitude.}
\begin{questions}
\item $g_0=\num{6,67e-11}\times\dfrac{\num{5,97e24}}{(\num{6,37e6})^2}=\res{\qty{9,81}{N.kg^{-1}}}$.
\item $P=m\,g_0=70\times9{,}81=\res{\qty{687}{N}}$.
\item $g=\num{6,67e-11}\times\dfrac{\num{5,97e24}}{(\num{6,77e6})^2}=\res{\qty{8,69}{N.kg^{-1}}}$, d'où $F=70\times8{,}69=\res{\qty{608}{N}}$, soit \qty{89}{\percent} du poids au sol.
\item \res{\text{Non}} : la Terre l'attire presque autant qu'au sol. Il « flotte » parce que la station et lui tombent ensemble autour de
la Terre, avec le même mouvement.
\end{questions}
\piege{Remplacer $d$ par l'altitude $h=\qty{400}{km}$ : on trouverait un champ gigantesque. La distance se mesure toujours entre les
\emph{centres} des corps.}
\end{corrige}
\end{blocLivre}
\begin{blocLivre}{cor-11-5}
\begin{corrige}{5}{Lire une carte de champ}
\demarche{Les lignes de champ partent des charges positives et arrivent sur les charges négatives. Le champ est tangent à la ligne,
dans son sens ; la force sur une charge $q<0$ est de sens opposé au champ.}
\begin{questions}
\item Les lignes partent de A : charge \res{\text{positive}} ; elles arrivent en B : charge \res{\text{négative}}.
\item En chaque point, on trace un vecteur tangent à la ligne qui y passe, orienté comme elle (en vert ci-dessous).
\item $\vec{F}=q\,\vec{E}$ avec $q<0$ : même direction que $\vec{E}$, \res{\text{sens opposé}} (en orange).
\item Les lignes sont plus \res{\text{serrées}} en M qu'en N : le champ est plus intense en \res{\text{M}}.
\end{questions}
\schema[corrige]{11-ex05}
\retiens{Une ligne de champ ne se coupe jamais avec une autre : en chaque point, le champ n'a qu'une seule direction.}
\end{corrige}
\end{blocLivre}
\begin{blocLivre}{cor-11-6}
\begin{corrige}{6}{Le champ de gravitation en Python}
\demarche{Le champ varie comme $\dfrac{1}{d^2}$, avec $d=R+h$ distance au centre de la Terre.}
\begin{questions}
\item La loi de Newton fait intervenir la distance au \res{\text{centre}} de la Terre, c'est-à-dire $R+h$.
\item $g=\num{6,67e-11}\times\dfrac{\num{5,97e24}}{(\num{1,637e7})^2}\approx\res{\qty{1,49}{N.kg^{-1}}}$.
\item $g$ est divisé par 4 quand la distance au centre double : $R+h=2R$, donc $h=R_\mathrm{T}=\res{\qty{6370}{km}}$.
\item Il suffit de remplacer la première ligne par \code{G, M, R = 6.67e-11, 7.35e22, 1.74e6} et de ne garder que \code{h = 0} :
le programme affiche \code{1.62 N/kg}.
\end{questions}
\end{corrige}
\end{blocLivre}
\begin{blocLivre}{cor-11-7}
\begin{corrige}{7}{L'atome d'hydrogène}
\demarche{Calculer les deux forces avec la même distance $r$, puis leur rapport : les $r^2$ se simplifient.}
\begin{questions}
\item $F_\mathrm{e}=k\,\dfrac{e^2}{r^2}=\num{9,0e9}\times\dfrac{(\num{1,60e-19})^2}{(\num{5,3e-11})^2}=\res{\qty{8,2e-8}{N}}$ ; charges opposées : force \res{\text{attractive}}.
\item $F_\mathrm{g}=G\,\dfrac{m_\mathrm{e}\,m_\mathrm{p}}{r^2}=\num{6,67e-11}\times\dfrac{\num{9,11e-31}\times\num{1,67e-27}}{(\num{5,3e-11})^2}=\res{\qty{3,6e-47}{N}}$.
\item $\dfrac{F_\mathrm{e}}{F_\mathrm{g}}\approx\res{\num{2,3e39}}$ : la cohésion de l'atome est assurée par l'interaction \res{\text{électrostatique}}.
\item Les astres sont globalement \res{\text{neutres}} : les forces électriques se compensent, alors que la gravitation, toujours
attractive, s'ajoute pour toutes les masses.
\end{questions}
\end{corrige}
\end{blocLivre}
\begin{blocLivre}{cor-11-8}
\begin{corrige}{8}{Deux lois si semblables}
\demarche{Comparer terme à terme les deux expressions vectorielles.}
\begin{questions}
\item $\vec{F}_{\mathrm{A}/\mathrm{B}}=k\,\dfrac{q_\mathrm{A}\,q_\mathrm{B}}{d^2}\,\vec{u}_{\mathrm{AB}}$ et $\vec{F}_{\mathrm{A}/\mathrm{B}}=-G\,\dfrac{m_\mathrm{A}\,m_\mathrm{B}}{d^2}\,\vec{u}_{\mathrm{AB}}$.
\item Les deux forces sont \res{\text{proportionnelles}} au produit des grandeurs (charges ou masses) et à $\dfrac{1}{d^2}$ ; elles agissent à
distance, le long de la droite qui joint les deux corps.
\item La gravitation est toujours \res{\text{attractive}}, l'interaction électrostatique peut être attractive ou répulsive ; la constante
$k$ est environ $10^{20}$ fois plus grande que $G$.
\item Avec $q_\mathrm{A}\,q_\mathrm{B}>0$, $\vec{F}_{\mathrm{A}/\mathrm{B}}$ a le sens de $\vec{u}_{\mathrm{AB}}$ : B est repoussé. Les masses étant positives, il faut le signe
« $-$ » pour que la force soit dirigée vers A : la gravitation est \res{\text{attractive}}.
\end{questions}
\end{corrige}
\end{blocLivre}
\begin{blocLivre}{cor-11-9}
\begin{corrige}{9}{Un électron dans un condensateur}
\demarche{Champ uniforme : $E=\dfrac{U}{d}$, de la plaque $+$ vers la plaque $-$ ; force $\vec{F}=q\,\vec{E}$ avec $q=-e$.}
\begin{questions}
\item $E=\dfrac{\num{1,5e3}}{\num{3,0e-2}}=\res{\qty{5,0e4}{V.m^{-1}}}$ ; lignes de champ parallèles, de la plaque supérieure vers la plaque inférieure.
\item $F=e\,E=\num{1,60e-19}\times\num{5,0e4}=\res{\qty{8,0e-15}{N}}$, verticale, vers la plaque \res{\text{positive}}.
\item $P=m_\mathrm{e}\,g=\num{9,11e-31}\times9{,}81=\res{\qty{8,9e-30}{N}}$ ; $\dfrac{F}{P}\approx\num{9,0e14}$ : le poids est \res{\text{négligeable}}.
\item Le proton, de charge $+e$, subirait une force de même valeur mais \res{\text{de sens opposé}}, vers la plaque négative.
\end{questions}
\schema[corrige]{11-ex09}
\piege{Écrire $\vec{F}=e\,\vec{E}$ pour un électron : sa charge vaut $-e$, la force est donc opposée au champ.}
\end{corrige}
\end{blocLivre}
\begin{blocLivre}{cor-11-10}
\begin{corrige}{10}{Le pendule électrostatique}
\demarche{Raisonner sur les électrons, seules charges mobiles dans un conducteur : ils fuient la règle négative et peuvent passer
de la règle à la boule au contact.}
\begin{questions}
\item Par \res{\text{influence}}, les électrons de la boule sont repoussés vers la face éloignée ; la face proche devient positive. Plus
proche de la règle, elle est plus attirée que la face négative n'est repoussée : la boule est \res{\text{attirée}}.
\item Au contact, des \res{\text{électrons}} passent de la règle sur la boule.
\item La boule est maintenant chargée négativement, comme la règle : elles se \res{\text{repoussent}}.
\item Les signes seraient inversés : attraction par influence, puis la boule céderait des électrons à la règle, deviendrait
positive et serait \res{\text{repoussée}}.
\end{questions}
\schema[corrige]{11-ex10}
\end{corrige}
\end{blocLivre}
\begin{blocLivre}{cor-11-11}
\begin{corrige}{11}{Entre la Terre et la Lune}
\demarche{En P, les champs ont même valeur et des sens opposés ; on égale leurs valeurs, puis on prend la racine carrée.}
\begin{questions}
\item Les deux champs ont la \res{\text{même valeur}} et des \res{\text{sens opposés}} : le champ total est nul (schéma ci-dessous).
\item $G\,\dfrac{M_\mathrm{T}}{x^2}=G\,\dfrac{M_\mathrm{L}}{(d_\mathrm{TL}-x)^2}$ : on simplifie par $G$.
\item $\dfrac{x^2}{(d_\mathrm{TL}-x)^2}=\dfrac{M_\mathrm{T}}{M_\mathrm{L}}=81{,}2$, donc $\dfrac{x}{d_\mathrm{TL}-x}=\sqrt{81{,}2}=9{,}01$.\par\smallskip
Ainsi $x=9{,}01\,(d_\mathrm{TL}-x)$, d'où $10{,}01\,x=9{,}01\,d_\mathrm{TL}$ :
\[x=\frac{9{,}01}{10{,}01}\times\num{3,84e8}=\res{\qty{3,46e8}{m}}\approx\res{0{,}90\,d_\mathrm{TL}}.\]
\item La Lune est 81 fois moins massive : son champ n'égale celui de la Terre que très près d'elle ($\sqrt{81}=9$ fois plus près).
\end{questions}
\schema[corrige]{11-ex11-cor}
\retiens{Pour résoudre $\dfrac{a}{x^2}=\dfrac{b}{(d-x)^2}$, on prend la racine carrée des deux membres : l'équation devient du premier degré.}
\end{corrige}
\end{blocLivre}
\begin{blocLivre}{cor-11-12}
\begin{corrige}{12}{La photocopieuse}
\demarche{Masse $m=\rho\,V$ avec $r=\qty{5,0}{\micro m}$ ; poids $P=m\,g$ ; force électrique $F=|q|\,E$ ; puis on compare.}
\begin{questions}
\item $n=\dfrac{\num{2,0e-14}}{\num{1,60e-19}}=\res{\num{1,25e5}}$ charges élémentaires.
\item $V=\dfrac{4}{3}\,\pi\,(\num{5,0e-6})^3=\qty{5,24e-16}{m^3}$, $m=\num{1,2e3}\times\num{5,24e-16}=\res{\qty{6,3e-13}{kg}}$, $P=\res{\qty{6,2e-12}{N}}$.
\item $F=\num{2,0e-14}\times\num{1,0e6}=\res{\qty{2,0e-8}{N}}$.
\item $\dfrac{F}{P}=\dfrac{\num{2,0e-8}}{\num{6,2e-12}}\approx\res{\num{3,2e3}}$ : la force électrique est plus de trois mille fois plus grande que le poids. Le
toner est donc arraché vers les zones chargées du tambour, quelle que soit leur orientation : l'image se forme fidèlement.
\end{questions}
\piege{Prendre le diamètre au lieu du rayon dans $V=\dfrac{4}{3}\pi r^3$ : le volume serait 8 fois trop grand.}
\end{corrige}
\end{blocLivre}
\begin{blocLivre}{cor-11-13}
\begin{corrige}{13}{Un satellite géostationnaire}
\demarche{$d=R_\mathrm{T}+h$, puis $g=G\,\dfrac{M_\mathrm{T}}{d^2}$ et $F=m\,g$. Avec $g_0=G\,\dfrac{M_\mathrm{T}}{R_\mathrm{T}^2}$, le rapport donne $g=g_0\,\dfrac{R_\mathrm{T}^2}{(R_\mathrm{T}+h)^2}$.}
\begin{questions}
\item $d=6370+35786=\res{\qty{42156}{km}}=\qty{4,2156e7}{m}$.
\item $g=\num{6,67e-11}\times\dfrac{\num{5,97e24}}{(\num{4,2156e7})^2}=\res{\qty{0,224}{N.kg^{-1}}}$.
\item $F=3000\times0{,}224=\res{\qty{672}{N}}$, contre $P_0=3000\times9{,}81=\qty{2,94e4}{N}$ au sol : environ \res{\text{44 fois}} moins.
\item $g=g_0\left(\dfrac{R_\mathrm{T}}{R_\mathrm{T}+h}\right)^2=9{,}81\times\left(\dfrac{6370}{42156}\right)^2=\res{\qty{0,224}{N.kg^{-1}}}$ : on retrouve bien le résultat.
\end{questions}
\retiens{Le rapport $\dfrac{g}{g_0}$ ne dépend que des distances : on peut calculer un champ en altitude sans connaître ni $G$ ni la masse de la Terre.}
\end{corrige}
\end{blocLivre}
\begin{blocLivre}{cor-11-14}
\begin{corrige}{14}{L'expérience de Millikan}
\demarche{Gouttelette immobile : les forces se compensent, $|q|\,E=m\,g$. Le champ entre les plaques vaut $E=\dfrac{U}{d}$.}
\begin{questions}
\item Le \res{\text{poids}} $\vec{P}$, vertical vers le bas, et la \res{\text{force électrostatique}} $\vec{F}=q\,\vec{E}$, verticale vers le haut.
\item La gouttelette est négative : la force est opposée au champ. Le champ est donc vers le bas : il va de la plaque positive
vers la plaque négative, la plaque \res{\text{supérieure}} est positive.
\item $E=\dfrac{600}{\num{1,0e-2}}=\res{\qty{6,0e4}{V.m^{-1}}}$ ; $|q|=\dfrac{m\,g}{E}=\dfrac{\num{2,94e-15}\times9{,}81}{\num{6,0e4}}=\res{\qty{4,81e-19}{C}}$.
\item $\dfrac{|q|}{e}=\dfrac{\num{4,81e-19}}{\num{1,60e-19}}\approx\res{3{,}0}$ : la gouttelette porte \res{3} électrons en excès.
\end{questions}
\schema[corrige]{11-ex14}
\end{corrige}
\end{blocLivre}
\begin{blocLivre}{cor-11-15}
\begin{corrige}{15}{Peser la Terre}
\demarche{Au sol, $g=G\,\dfrac{M_\mathrm{T}}{R_\mathrm{T}^2}$ : on isole $M_\mathrm{T}$. La masse volumique vaut $\rho=\dfrac{M}{V}$.}
\begin{questions}
\item $\res{M_\mathrm{T}=\dfrac{g\,R_\mathrm{T}^2}{G}}$.
\item $M_\mathrm{T}=\dfrac{9{,}81\times(\num{6,37e6})^2}{\num{6,67e-11}}=\res{\qty{5,97e24}{kg}}$.
\item $V=\dfrac{4}{3}\,\pi\,(\num{6,37e6})^3=\qty{1,08e21}{m^3}$, d'où $\rho=\res{\qty{5,51e3}{kg.m^{-3}}}$ : environ deux fois celle des roches de surface.
\item L'intérieur de la Terre est formé de matériaux \res{\text{beaucoup plus denses}} que la croûte : c'est le noyau, riche en fer.
\end{questions}
\end{corrige}
\end{blocLivre}
